\section{Generace audio výstupu}
\label{pp:audio}

V podkapitole \ref{sec:dac} bylo zmíněno, že pro generaci audio signálu z výstupu APU je využívána
pulsně šířková modulace (PWM) pro aproximaci audio signálu v digitální doméně. Naše řešení kromě PWM
také využívá technologii DMA pro synchronní a stabilní přenos vzorku podle vzorkovací frekvence.

\subsection{Konfigurace PWM}
\label{sub:pwm_configuration}

Pro dosažení kvalitního a stabilního zvukového výstupu je nezbytné zajistit, aby byly jednotlivé
vzorky odesílány do PWM generátoru přesně v rytmu vzorkovací frekvence $f_s$. Jakákoliv softwarová
nepravidelnost (tzv. \textit{jitter}) by se projevila jako slyšitelné zkreslení (trhaní zvuku,
zasekávání atd.). Z tohoto důvodu byla zvolena plně hardwarová cesta využívající řetězení kanálů
DMA, díky čemuž procesor věnuje tomu minimální svůj čas a je menší šance desynchronizace.

Nejprvé je nutné odvodit parametry PWM z frekvence systémových hodin $f_{\cmrtext{clk}}$. Pro
dosažení optimální filtrace v analogové části obvodu byla zvolena
nosná frekvence PWM $f_{\cmrtext{PWM}} = \SI{250}{\kilo\hertz}$. Vzorkovací frekvence emulátoru je
pevně dána na $f_s = \SI{31.25}{\kilo\hertz}$. Z toho vyplývá, že na každý jeden audio vzorek
připadá přesně 8 celých PWM cyklů:

$$
\frac{f_{\text{PWM}}}{f_s} = \frac{250000}{31250} = 8
$$

Nosná Frekvence PWM byla zvolena jako celočíselný násobek vzorkovací frekvence , aby bylo zajištěno
deterministické a periodické časování mezi aktualizací střídy PWM a jednotlivými vzorky audio
signálu. Tím je zaručeno, že každý vzorek je reprezentován konstantním počtem period PWM, což
eliminuje časovou nekonzistenci při přepisu hodnot (jitter) a zabraňuje vzniku nízkofrekvenčních
artefaktů způsobených nesynchronním vzorkováním. V případě neceločíselného poměru by docházelo k
postupnému fázovému driftu mezi PWM signálem a generováním vzorků, což by vedlo k periodickému
kolísání efektivní šířky pulzu a degradaci kvality rekonstruovaného analogového signálu.

Klíčovým parametrem PWM je hodnota stropu čítače (TOP), která definuje maximální rozlišení (bitovou
hloubku) výstupního signálu, závislé na počtu možných stavu stropu čítače TOP:

$$
\cmrtext{bitová hloubka} = \log_2(\cmrtext{TOP} + 1)
$$

Vztah pro výpočet frekvence PWM v mikrokontroléru RP2350 je definován jako:

$$
f_\cmrtext{PWM} = \frac{f_\cmrtext{clk}}{\cmrtext{clkdiv} \cdot (\cmrtext{TOP} + 1)}
$$

Jelikož nosní frekvenci PWM $f_\cmrtext{PWM}$ máme vybranou a taktovací frekvenci RP2350
$f_\cmrtext{clk}$ známe, potřebujeme vypočítat optimální hodnotu pro strop čítače TOP. Po
přeskládaní rovnice a dosazení hodnot získáváme:

$$
\cmrtext{TOP} = \frac{f_\cmrtext{clk}}{f_\cmrtext{PWM}} - 1 = \frac{300 \cdot 10^6}{250 \cdot 10^3} - 1 = 1200 - 1 = 1199
$$

Rozlišení PWM výstupu je tedy 1200 úrovní (včetně 0). Bitovou hloubku získáme hodnotu přibližně 10,2
bitů ($\log_2(1200) \approx 10,2$). Ačkoliv moderní audio systémy využívají 16bitové nebo 24bitové
rozlišení, pro reprodukci 8bitového zvuku z konzole NES je 10bitová hloubka více než dostatečná.

Jelikož APU generuje 16bitové vzorky (v rozsahu 0 až 65535), je nutné je před odesláním do
komparátoru přemapovat na rozsah čítače (0 až 1199):

$$
\cmrtext{hodnota} = \left\lfloor{\frac{\cmrtext{vzorek} \cdot \cmrtext{TOP}}{65535}}\right\rfloor
$$

Funkce \texttt{init\_pwm()} ve výpisu \ref{lst:pwm_init} provádí konfiguraci a inicializaci PWM
podle výše uvedených informací.

\begin{listing}[htbp]
\begin{minted}{c}
/* Konvertace vzorku na hodnotu čítače. Přetypování na uint32_t
 * je zapotřebí, pokud výsledek násobení bude větší než 65535. */
#define PWM_VALUE(smp) \
    ((uint16_t) ((((uint32_t) (smp) * g_pwm_wrap) / 65535)))

#define PWM_FREQ 250000

static unsigned g_pwm_slice; /* číslo slicu */
static uint16_t g_pwm_wrap; /* TOP */

void init_pwm(void) {
    /* Nastavení pinu jako výstup PWM */
    gpio_set_function(APU_OUT_PIN, GPIO_FUNC_PWM);

    g_pwm_slice = pwm_gpio_to_slice_num(APU_OUT_PIN);

    /* Výpočet TOP (wrap v terminologii RP2350) */
    unsigned sys_clk = clock_get_hz(clk_sys);
    g_pwm_wrap = sys_clk / PWM_FREQ - 1;
    /* Inicializace PWM */
    pwm_set_wrap(g_pwm_slice, g_pwm_wrap);
    pwm_set_clkdiv(g_pwm_slice, 1.f); /* dělička je 1 */
    pwm_set_enabled(g_pwm_slice, true);
}
\end{minted}
\caption{Inicializace a konfigurace PWM.}
\label{lst:pwm_init}
\end{listing}

\subsection{Architektura DMA přenosu}

Aby byl zajištěn nepřerušený tok dat, je využit mechanismus dvojitého bufferování, známý také jako
\textbf{ping-pong buffer}\cite{pingpong-buf}. Paměťový prostor pro vzorky je rozdělen do dvou polí.
Zatímco emulátor plní novými daty jeden buffer, DMA řadič čte data z druhého bufferu a posílá
je do PWM komparátoru.

\newpage

K řízení tohoto procesu jsou zapotřebí dva DMA kanály (obr. \ref{fig:dma_channels}):

\begin{enumerate}
    \item \textbf{Audio kanál} -- tento kanál je zodpovědný za přenos samotných vzorků (16bitových
        hodnot) z aktuálního bufferu přímo do hardwarového registru PWM komparátoru. Přenos jednoho
        vzorku je spouštěn speciálním DMA časovačem, který je nastaven přesně na vzorkovací
        frekvenci $f_s = \SI{31.25}{\kilo\hertz}$.
    \item \textbf{Řídicí kanál} -- jakmile audio kanál vyprázdní aktuální buffer, automaticky spustí
        tento řídicí kanál. Jeho jediným úkolem je přepsat zdrojovou adresu audio kanálu na adresu
        druhého bufferu. Následně řídicí kanál opětovně spustí audio kanál, čímž se cyklus uzavírá
        do nekonečné smyčky.
\end{enumerate}

Tento řetězený přístup zaručuje, že nedochází k žádnému softwarovému zpoždění při přepínání bufferů.
Alternativným řešením by bylo zaregistrovat přerušení, které by se vyvolalo když jediný DMA audio
kanál dokončí přenos dat. Ve vývoji se ale ukázalo, že to naopak způsobuje zpomalení procesoru, čímž
je ovlivněna plynulost běhu emulaci, a také vyvolává trhání zvuku.

\begin{figure}[htbp]
    \centering
    \includesvg[width=0.75\textwidth]{images/dma_channels.svg}
    \caption{Blokové schéma principu dvou DMA kanálu: audio kanál a řídicí kanál.}
    \label{fig:dma_channels}
\end{figure}

\subsection{Konfigurace DMA}

Teoretický koncept ping-pong bufferu a řetězení kanálů je v realizován pomocí sady funkcí pico-sdk.

Při inicializaci DMA, je nutné nastavit frekvenci DMA časovače $f_\cmrtext{DMA}$ na potřebnou.
Obecný vztah pro $f_\cmrtext{DMA}$ platí:

$$
f_\cmrtext{DMA} = f_\cmrtext{clk} \frac{n}{d}
$$

V našem případě chceme docílit toho, aby frekvence časovače DMA se rovnala vzorkovací frekvenci,
takže potřebujeme se zbavit členu $f_\cmrtext{clk}$ z rovnice. Když dosadíme $n = f_s$ a $d =
f_\cmrtext{DMA}$, dostaneme:

$$
f_\cmrtext{DMA} = f_\cmrtext{clk} \frac{n}{d} = 300 \cdot 10^6 \cdot \frac{3,125 \cdot 10^4}{300
\cdot 10^6} = \SI{31,25}{kHz}
$$

Knihovna pico-sdk poskytuje funkci \texttt{dma\_timer\_set\_fraction(timer, n, d)} pro nastavení
frekvence DMA časovače. Jelikož je ale taktovací frekvence jádra 300 MHz, tato hodnota přesahuje
maximální limit 16bitového dělitele. Problém byl vyřešen zkrácením zlomku na poměr
$\frac{1}{9600}$, což přesně odpovídá požadovaným 31,25 kHz.

\subsection{Generace vzorků}

Zápis vzorků do správného bufferu zajišťuje funkce \texttt{audio\_queue()} (výpis
\ref{lst:dma_queue}). Pokaždé když APU vygeneruje nový vzorek, je tato funkce zavolána. Jakmile
dojde k naplnění aktuálního bufferu, musí se přejít na další buffer.

Tato funkce také zajišťuje synchronizaci v okamžiku, když CPU zaplní buffer a chce přepnout se na
další: CPU zkontroluje hardwarový registr \texttt{read\_addr} DMA řadiče. Pokud DMA stále čte z
druhého bufferu, znamená to, že CPU emuluje systém rychleji, než probíhá reálné přehrávání zvuku.
Procesor se proto zablokuje ve smyčce \texttt{tight\_loop\_counter()}, dokud se buffer neuvolní.
Tímto způsobem emulátor se také synchronizuje se zvukem a udržuje stabilní rychlost běhu bez
nutností využití softwarových časovačů.

\begin{listing}[htbp]
\begin{minted}{c}
void audio_queue(ApuSample smp) {
    /* Přemapování 16bitového vzorku na rozlišení PWM a zápis do bufferu */
    g_bufs[g_curr_buf_idx][g_buf_pos++] = PWM_VALUE(smp);

    /* Pokud je aktuální buffer plný, připravíme se na přepnutí */
    if (g_buf_pos >= AUDIO_BUFFER_SIZE) {
        unsigned next_buf_idx = (g_curr_buf_idx + 1) & (RING_SIZE - 1);

        for (;;) {
            /* Získání aktuální adresy, ze které DMA právě čte */
            uintptr_t dma_addr = dma_hw->ch[g_audio_chan].read_addr;
            uintptr_t next_buf_start = (uintptr_t) g_bufs[next_buf_idx];
            uintptr_t next_buf_end = next_buf_start + sizeof(g_bufs[0]);

            /* Pokud DMA už z dalšího bufferu nečte, můžeme smyčku opustit */
            if (dma_addr < next_buf_start || dma_addr >= next_buf_end)
                break;

            /* Emulátor běží příliš rychle, CPU musí počkat na hardware audia */
            tight_loop_contents();
        }

        /* Přepnutí indexů pro další dávku vzorků */
        g_curr_buf_idx = next_buf_idx;
        g_buf_pos = 0;
    }
}
\end{minted}
\caption{Ukládání vzorku do fronty (volného bufferu).}
\label{lst:dma_queue}
\end{listing}
